% !TEX program = xelatex
% Evidence-first market research report template.
\documentclass[11pt,letterpaper]{report}
\usepackage{market_research}

\newcommand{\reporttitle}{[Market name]}
\newcommand{\reportsubtitle}{Evidence-first market assessment}
\newcommand{\retrievalcutoff}{[YYYY-MM-DD]}
\newcommand{\reportauthor}{[Prepared by]}
\newcommand{\reportclassification}{[Public / Internal / Confidential]}

\hypersetup{
  pdftitle={\reporttitle{} --- \reportsubtitle{}},
  pdfauthor={\reportauthor{}},
  pdfkeywords={market definition, evidence ledger, scenario analysis}
}

\begin{document}

\makemarketreporttitle
  {\reporttitle}
  {\reportsubtitle}
  {\retrievalcutoff}
  {\reportauthor}

\pagenumbering{roman}

\begin{calloutbox}[Use and limitations]
This report distinguishes observed facts, estimates, calculations, scenarios,
opinions, and recommendations. Claim IDs map to the claims ledger; source IDs
map to the evidence ledger. Scenario ranges are conditional and are not
confidence intervals unless a separately documented statistical model supports
that interpretation. This report is not investment, legal, or financial advice.
\end{calloutbox}

\textbf{Classification:} \reportclassification

\tableofcontents
\clearpage
\pagenumbering{arabic}

\chapter{Executive Synopsis}

\section{Decision Context}
[State the decision, audience, retrieval cutoff, and material scope constraints.]

\section{Findings}
\begin{evidencebox}[Observed findings]
\begin{itemize}
  \item [Finding with \claimref{C-001} and \sourceref{S-001}.]
  \item [Finding with explicit geography, period, unit, and revision status.]
\end{itemize}
\end{evidencebox}

\section{Scenario Range}
\begin{calculationbox}[Conditional range]
[State the range across named scenarios and methods. Include calculation IDs,
currency, base year, price basis, period, and denominator.]
\end{calculationbox}

\begin{assumptionbox}[Key assumptions]
[List the assumptions that most influence the result and their
\assumptionref{A-...} identifiers.]
\end{assumptionbox}

\begin{limitationbox}[Material limitations]
[List unresolved source conflicts, missing evidence, revisions, sample
limitations, taxonomy breaks, and model limitations.]
\end{limitationbox}

\chapter{Scope, Definitions, and Method}

\section{Market Definition}
\begin{calloutbox}[Formal definition]
[Define the product or service, customer, geography, transaction type,
measurement basis, and period.]
\end{calloutbox}

\subsection{Included}
\begin{itemize}
  \item [Included product, customer, or channel.]
\end{itemize}

\subsection{Excluded}
\begin{itemize}
  \item [Excluded category and reason.]
\end{itemize}

\section{Measurement Contract}
\begin{longtable}{@{}p{0.28\textwidth}p{0.64\textwidth}@{}}
\toprule
Field & Definition \\
\midrule
\endhead
Geography & [Boundary and treatment of imports/exports] \\
Historical period & [Start--end] \\
Forecast period & [Start--end] \\
Currency and base year & [Currency; nominal/real/current/constant basis] \\
Measure type & [Stock/flow/count/share/rate/price/index] \\
Unit and denominator & [Explicit unit and denominator ID] \\
Industry taxonomy & [NAICS/NACE/ISIC or other taxonomy and version] \\
Product taxonomy & [NAPCS/CPA/PRODCOM/CPC or other taxonomy and version] \\
Retrieval cutoff & \retrievalcutoff \\
\bottomrule
\end{longtable}

\section{Evidence Hierarchy and Conflicts}
[Describe the source hierarchy, exact claim--source mapping, snapshot/archive
policy, revision treatment, and how conflicting sources were reconciled or
retained as a range.]

\chapter{Market Size and Reconciliation}

\section{Top-Down Method}
[List disjoint components, coverage keys, sources, conversions, and exclusions.
Do not sum gross output, revenue, trade, and end-customer spend as though they
were the same denominator.]

\section{Bottom-Up Method}
[Show customer or unit counts, annual quantities, price assumptions,
addressable fractions, sources, and formulas.]

\section{TAM, SAM, and SOM Scenarios}
\begin{markettable}{Conditional TAM/SAM/SOM scenarios}
\begin{tabular}{@{}lrrrl@{}}
\toprule
Scenario & TAM & SAM & SOM & Calculation \\
\midrule
Downside & [value] & [value] & [value] & \calcref{CALC-...} \\
\tablerowcolor
Base & [value] & [value] & [value] & \calcref{CALC-...} \\
Upside & [value] & [value] & [value] & \calcref{CALC-...} \\
\bottomrule
\end{tabular}
\end{markettable}

\begin{assumptionbox}[Interpretation]
TAM, SAM, and SOM are conditional scenario outputs. They are not a single
asserted market truth. Explain serviceability filters, capture constraints,
time horizon, and why the scenarios differ.
\end{assumptionbox}

\section{Method Reconciliation}
[Report the absolute gap and gap as a percentage of the methods' midpoint.
Investigate inconsistent boundaries, denominators, taxes, channels, time
periods, currencies, and double counting before averaging estimates.]

\section{Sensitivity}
[Show how results change when the principal assumptions vary. State switching
values at which a decision would change.]

\chapter{Demand and Customer Evidence}

\section{Observed Demand}
[Report administrative, transactional, official statistical, or filing-based
evidence with exact definitions and source IDs.]

\section{Survey Evidence}
\begin{longtable}{@{}p{0.30\textwidth}p{0.62\textwidth}@{}}
\toprule
Disclosure & Detail \\
\midrule
\endhead
Sponsor and fieldwork organization & [Names] \\
Target population and sampling frame & [Definition] \\
Probability or non-probability sample & [Method] \\
Recruitment and mode & [Method, languages, dates] \\
Unweighted sample and subgroups & [Counts] \\
Weighting and design effect & [Variables, benchmarks, effect] \\
Response or participation & [Definition and rate, if available] \\
Question wording & [Instrument location] \\
Precision & [Appropriate measure and assumptions] \\
Limitations & [Coverage, nonresponse, measurement, model error] \\
\bottomrule
\end{longtable}

\section{Interview Evidence}
[Describe recruitment, consent, dates, roles represented, analysis/coding
method, saturation limitations, and privacy controls. Present themes as
qualitative evidence; do not attach population prevalence to them.]

\chapter{Competitive Landscape}

\section{Relevant Product and Geographic Scope}
[Explain substitutability, customer perspective, non-price competition,
channels, geography, dynamic change, and alternate plausible definitions.
Descriptive market research does not determine a legal antitrust market.]

\section{Competitor-Feature Matrix}
[Insert the validated, evidence-linked matrix. Use ``unknown'' when public
evidence is absent. Keep product edition, geography, and as-of date constant.]

\section{Shares and Concentration}
[State the share metric and denominator: revenue, units, capacity, active
users, or another defensible measure. Explain residual/unknown participants and
source coverage. If HHI is shown, state the exact market definition and avoid
turning screening indicators into legal conclusions.]

\section{Competitive Dynamics}
[Separate observed actions from interpretation. Cite filings, regulator
records, official registries, or attributable public company material.]

\chapter{Forecast Scenarios}

\section{Historical Series and Revisions}
[State source series, frequency, seasonal adjustment, vintage/retrieval date,
revisions, transformations, and any taxonomy or methodology breaks.]

\section{Scenario Definitions}
\begin{longtable}{@{}p{0.18\textwidth}p{0.34\textwidth}p{0.38\textwidth}@{}}
\toprule
Scenario & Rate path or drivers & Assumptions and evidence \\
\midrule
\endhead
Downside & [Path] & [Assumption and source IDs] \\
Base & [Path] & [Assumption and source IDs] \\
Upside & [Path] & [Assumption and source IDs] \\
\bottomrule
\end{longtable}

\section{Range and Sensitivity}
[Report the conditional range by year, endpoint sensitivity, breakpoints, and
which input dominates. Do not assign scenario probabilities without a
validated probabilistic basis.]

\chapter{Regulation, Risks, and Uncertainty}

\section{Regulatory Record}
[Distinguish enacted, effective, proposed, consulted, stayed, and repealed
rules. Record jurisdiction, authority, publication date, effective date, and
primary legal source.]

\section{Risk Register}
[For each risk, separate evidence, likelihood judgment, impact mechanism,
early indicators, and mitigation. Avoid false numerical precision.]

\chapter{Implications and Options}

\section{Implications}
[Link each implication to findings and state what would disconfirm it.]

\section{Options}
[Compare options against explicit objectives, dependencies, costs, risks, and
decision thresholds.]

\begin{recommendationbox}[Recommendations]
[State recommendations as judgments, not sourced facts. Include dependencies,
owners, timing, and evidence that would trigger reconsideration.]
\end{recommendationbox}

\chapter{Limitations}
[Consolidate evidence gaps, source conflicts, potential double counting,
currency/base-year conversions, revision exposure, survey error, interview
limits, competitive-intelligence limits, and forecast/model uncertainty.]

\appendix

\chapter{Evidence Ledger}
[Render the complete source ledger: source ID, title, publisher, URL or
persistent identifier, publication and retrieval dates, geography, unit,
currency/base, taxonomy, revision status, method/sample, limitations, terms,
and archive path.]

\chapter{Claims Ledger}
[Render claim ID, exact claim text, statement type, source IDs, calculation and
assumption IDs, report location, scope metadata, revision status, and
confidence label.]

\chapter{Calculations and Assumptions}
[Include formulas, raw inputs, conversions, coverage keys, denominator IDs,
scenario paths, sensitivity values, and reconciliation.]

\chapter{Research Ethics and Data Handling}
[Document consent, privacy notices, retention, access controls,
de-identification, lawful public-source collection, and deletion schedule.
Confirm that no deceptive collection, PII disclosure, or trade-secret
acquisition was used.]

\end{document}
